\section{Architektura}
\label{sec:architektura}

Kod jest pakietem Pythona \kod{moderator} w katalogu \kod{src/}. Każdy plik robi jedną rzecz, a terminal i aplikacja
webowa korzystają z jednej klasy \kod{Moderator} (plik \kod{service.py}), więc wynik jest taki sam niezależnie od
interfejsu.

\begin{table}[H]
\centering
\small
\begin{tabularx}{\linewidth}{@{}l X@{}}
\toprule
Plik & Zadanie \\
\midrule
\kod{config.py} & wczytanie \kod{config.yaml} i zmiennych środowiskowych (\kod{.env}) \\
\kod{llm.py} & klient Ollamy: czat, czat strumieniowy, obraz, embeddingi \\
\kod{schema.py} & schemat odpowiedzi (Pydantic) i jego walidacja \\
\kod{prompts.py} & prompt systemowy, przykłady few-shot, szablony RAG, odczytu obrazu i czatu \\
\kod{classify.py} & ocena komentarza (zero-shot, few-shot) z ponowieniem przy złym JSON-ie \\
\kod{rules.py}, \kod{retrieval.py}, \kod{rag.py} & podział regulaminu na paragrafy, BM25, embeddingi, RRF, ocena z cytatem \\
\kod{vision.py} & odczyt tekstu z obrazka modelem wizyjnym \\
\kod{baseline.py} & model bazowy \kod{unitary/toxic-bert} \\
\kod{batch.py}, \kod{report.py} & moderacja pliku, raport \kod{.txt}, wykres \kod{.svg} \\
\kod{chat.py} & sesja rozmowy o pliku z historią i komendami \\
\kod{evaluate.py} & metryki porównania wariantów \\
\kod{cli/} & polecenia terminala (Typer + Rich), podzielone na \kod{single}, \kod{files}, \kod{tools}, \kod{views} \\
\kod{web/} & aplikacja Gradio; każda zakładka w osobnym pliku w \kod{web/tabs/} \\
\bottomrule
\end{tabularx}
\caption{Moduły pakietu \kod{moderator}.}
\end{table}

Przepływ jednego komentarza w trybie RAG wygląda tak:
\begin{enumerate}
  \item (tylko obrazek) \kod{gemma3:4b} odczytuje tekst z obrazka,
  \item wyszukiwarka zwraca 3 najlepsze paragrafy regulaminu; do kontekstu zawsze dochodzi też \kod{§2.3}
        (dozwolona krytyka), żeby model miał podstawę do zostawienia wpisu,
  \item \kod{qwen2.5:7b} dostaje prompt systemowy, przykłady few-shot, paragrafy i komentarz, a odpowiada JSON-em
        zgodnym ze schematem,
  \item odpowiedź jest sprawdzana: zły JSON -- jedna ponowna próba, potem decyzja \kod{review}; paragraf, którego
        nie ma w regulaminie, albo usunięcie bez paragrafu -- decyzja zmieniana na \kod{review}.
\end{enumerate}

Dane są w całości syntetyczne: 40 komentarzy (23 po polsku, 17 po angielsku, 22 naruszające zasady) w
\kod{data/comments.csv} z etykietą, kategorią i właściwym paragrafem, fikcyjny regulamin forum ,,Kod po godzinach''
(14 paragrafów w \kod{data/rules/forum\_rules.md}) oraz 6 obrazków wygenerowanych skryptem
\kod{scripts/make\_sample\_images.py}. Komentarze toksyczne są celowo łagodne i nie zawierają prawdziwych wyzwisk wobec
grup.
